\chapter{Baby Universes}
\label{chap:baby}

\begin{quote}
\itshape
Universes are not born, they are extruded.
\end{quote}

\section{What forms after the bounce}

When matter reaches Cartasis density and the torsion repulsion overwhelms
gravity, the bouncing matter does not return to the parent's exterior (the
parent's horizon is now in its past, geometrically). Instead, the bounce
dynamics generate a new spacetime region---a baby universe---with its own time
direction and spatial extent.

From inside, the baby universe looks like a standard hot-Big-Bang cosmology:

\begin{itemize}
  \item it begins at maximum density (the bounce moment);
  \item it expands rapidly (torsion-driven, qualitatively similar to inflation);
  \item it cools as it expands;
  \item it transitions to standard radiation-dominated and matter-dominated phases;
  \item it eventually shows large-scale structure formation.
\end{itemize}

From outside (the parent's perspective), the baby is not visible as a spatial
structure. The parent sees a black hole at the location where the collapse
happened. The baby exists ``in the future of the bounce''---its spatial extent
is in directions orthogonal to the parent's interior geometry.

\section{The Tardis effect}

A consequence of the geometry: the parent black hole, viewed from outside, has a
specific Schwarzschild radius. The baby universe inside has its own cosmological
extent, which can be vastly larger than the parent's horizon would suggest.

For our universe, total mass--energy is $\sim 10^{53}\,\mathrm{kg}$. At Cartasis
density, this fits in $\sim 10^{3}\,\mathrm{m^3}$---less than half an Olympic
swimming pool. But from inside, we observe a universe $\sim 46$ billion
light-years in radius (the observable horizon).

This is not a contradiction. The two volumes are measured in different spacetime
regions. The ``small volume'' measures the matter at the bounce moment in parent
coordinates. The ``large volume'' measures the cosmological extent in baby
coordinates after expansion. They are not the same quantity.

This is also not unique to Einstein--Cartan. Even in standard general relativity,
the spatial volume inside a black hole's horizon grows with time and becomes
essentially unbounded~\citep{christodoulou2015}. The bounce framework
reinterprets this growth as cosmological expansion of the baby rather than
approach to a singularity.

\section{Inflation substitute}

The early baby universe undergoes rapid expansion because torsion repulsion is
at its maximum just after the bounce, when density is still very high. As the
baby expands and density drops, the torsion repulsion fades. Gravity reasserts
dominance, and expansion decelerates.

This phase has the right qualitative features to substitute for the inflationary
epoch of standard cosmology:

\begin{itemize}
  \item brief, very rapid expansion;
  \item ends naturally when conditions change (density drops below the Cartasis threshold);
  \item no separate inflaton field required;
  \item sets up homogeneous and approximately isotropic conditions for subsequent
    cosmological evolution.
\end{itemize}

Quantitative work is needed: derive the power spectrum of perturbations produced
by torsion-driven expansion and compare to inflation's predictions and to CMB
observations. This is the central observational test of the bounce framework
(see Chapters~\ref{chap:tests} and~\ref{chap:sims}).

\section{What the baby inherits from the parent}

Several quantities are transmitted across Cartasis.

\paragraph{Energy--momentum.} Conservation laws ensure the baby's total
energy--momentum equals what fell into the bounce. This sets the baby's
mass--energy budget.

\paragraph{Spin and angular momentum.} The parent's net angular momentum (if
rotating) is transmitted through the bounce. A rotating parent (Kerr-like)
produces a baby with a preferred axis---possibly the source of cosmological
anisotropies like the galaxy-rotation asymmetry observed in JWST
data~\citep{shamir2025}.

\paragraph{Chirality bias.} The bounce is CPT-even, so it neither flips charge
nor inherits a parent's asymmetry by relabeling. Instead the parity-violating
chiral-vortical filter at the rotating bounce \emph{generates} a matter bias in
the baby, of magnitude set by the bounce vorticity and sign set by its direction
(see Chapter~\ref{chap:chirality}).

\paragraph{Quantum entanglement.} Particles in the parent's near-horizon region
remain entangled with their Hawking-radiation partners. This entanglement
persists across the bounce, connecting the baby's interior with the parent's
exterior---see Chapter~\ref{chap:darksector}.

\paragraph{Conserved charges.} Baryon number, lepton number, and electric charge
(if conserved at trans-EW temperatures) flow through.

\medskip
\noindent What is \emph{not} transmitted:

\begin{itemize}
  \item structure (atoms, molecules, nuclei)---all dismantled at Cartasis density;
  \item specific particle identities---only conserved currents survive;
  \item information about the parent's ``before''---the bounce is the baby's past
    boundary.
\end{itemize}

\section{The persistent channel}

As discussed in Chapter~\ref{chap:cartasis}, the bounce is not a single event
but an ongoing interface. While the parent black hole continues to accrete
matter, that matter falls in, reaches Cartasis, and crosses into the baby. From
the baby's perspective, this means matter continues to be injected throughout
the parent's lifetime.

The injection geometry is non-trivial: new matter does not appear at a specific
spatial location in the baby's 3D space. It appears ``at the bounce membrane,''
which from the baby's perspective is the cosmic horizon at maximum lookback. So
injected matter manifests as gravitational influence from the cosmic boundary,
not as visible particles appearing in our 3D space. This is the basis for the
dark-matter interpretation in Chapter~\ref{chap:darksector}.

\begin{quote}
\small
\textbf{A note on inter-universe travel (a thought-experiment).}\enspace
A natural question---could one, granted every speculative tool that general
relativity and quantum theory permit (Alcubierre warp drives, traversable wormholes,
negative-energy/exotic matter), \emph{travel} to the parent universe, or visit a baby
universe forming in our sky? Posing it sharply exposes what the membranes really are,
and the answer is no in both directions, for the same reason. The obstruction is not
distance, which a warp drive defeats; it is that the membrane is not a wall in space
but a surface in \emph{time}---a bounce, an arrow-of-time boundary.

\emph{Outward (to the parent): forbidden.} Our Cartasis membrane is our own
Big-Bang surface---our past cosmological horizon. The parent's exterior lies
\emph{behind} that horizon and in our \emph{past}. A warp bubble can carry you
superluminally anywhere in your future light cone, but the parent is not in your
future; you cannot drive to your own Big Bang. Even the one genuine loophole the
framework already invokes---an ER\,$=$\,EPR / traversable-wormhole bridge for
information across the membrane (Chapter~\ref{chap:supraverse})---ferries only
pre-shared entanglement under a coupling you do not control, and it would open onto
the parent's bounce era, not its present; once the parent evaporates, the channel
pinches off entirely (Chapter~\ref{chap:supraverse}, Section~\ref{sec:recursion-depth}).

\emph{Inward (to a baby): you can seed one, not visit one.} A baby universe is a real
black hole in our sky; gravity will take you to it for free, and exotic-matter
shielding might, in this game, let you survive the bounce intact rather than being
thermalised at $\rhoc$. But surviving only means you arrive at the baby's
\emph{initial surface}: its time runs forward from its bounce, so the horizon delivers
you to its low-entropy Big Bang, not to some interior ``now.'' You would become part
of its first instant---a speck in its initial conditions (and, literally, part of the
inherited baryon debris its eventual observers measure as $\eta$,
Chapter~\ref{chap:chirality}). You can be an ancestor, never a tourist.

So warp drives, wormholes, and exotic matter all operate \emph{within} one
universe's thermodynamic arrow; none runs you \emph{backward across a low-entropy
initial surface}, because that is a direction in time and entropy, not a distance. You
may fall \emph{down} the recursion and seed a universe; you can never climb back
\emph{up} it. The membrane-as-Big-Bang structure is, in effect, the framework's own
chronology protection: it forbids the closed timelike curves that inter-generational
travel would otherwise allow.
\end{quote}

\section{Lifecycle}

A baby universe's lifetime is bounded by its parent black hole's Hawking
evaporation timescale. For a parent of mass $M$,
\[
  \tau_H \sim \frac{5120\,\pi\, G^2 M^3}{\hbar c^4}.
\]
For a stellar-mass parent, $\tau_H \sim 10^{67}$ years. For a supermassive
parent ($10^{9}\,\Msun$), $\tau_H \sim 10^{94}$ years. For our universe's mass
($\sim 10^{53}\,\mathrm{kg}$), $\tau_H \sim 10^{145}$ years.

These vastly exceed any internal cosmological timescale of the baby
($\sim 10^{10}$ years for matter--radiation transitions, $\sim 10^{100}$ years
for stellar evaporation). So the baby's internal physics plays out fully before
parent evaporation becomes relevant.

But the baby's lifetime is finite. At some point the parent's mass is mostly
radiated away, the bounce membrane shrinks, and the baby's connection to the
parent fades. What happens to the baby at that point is unclear---possibly
internal physics continues, possibly the baby's spacetime structure also
dissolves.

\section{Generation of black holes within the baby}

The baby's internal physics, if similar to ours, produces stellar collapse,
supernovae, and astrophysical black holes. Each such black hole, in the bounce
framework, is itself a potential progenitor of further baby universes---
grandchildren in the supraverse lineage.

Stellar-mass black holes ($\sim 10\,\Msun$) have Cartasis-density throat regions
that produce relatively small grandchildren. Supermassive black holes
($\sim 10^{9}\,\Msun$) produce larger grandchildren, potentially comparable to
our own universe in scale.

This cascading structure means each viable universe produces many descendants,
and the supraverse fills with universe-trees rooted in OG bounces (see
Chapter~\ref{chap:supraverse}).
